\documentclass[12pt]{article}
\usepackage[utf8]{inputenc}
\usepackage[T1]{fontenc}
\usepackage[margin=1in]{geometry}
\usepackage{setspace}
\usepackage{xcolor}
\usepackage[colorlinks=true, urlcolor=blue!60!black, linkcolor=blue!60!black]{hyperref}
\setlength{\parindent}{2em}
\setlength{\parskip}{1.25em}
\renewcommand{\baselinestretch}{1.0}
\begin{document}
\begin{center}
{\LARGE \textit{U.S. Department of Agriculture v. Moreno}, 413 U.S. 528 (1973)\par}
\vspace{0.8em}
{\large Case Brief --- Four-Page Edition\par}
\end{center}
\vspace{1em}

\section{Memory Jogger}

Congress cut off food stamps to any household containing an unrelated member, aiming at ``hippie communes''; the Court held 7--2 that the classification failed even rational-basis review, because a bare congressional desire to harm a politically unpopular group is not a legitimate governmental interest (\href{https://supreme.justia.com/cases/federal/us/413/528/}{Justia}).

\section{Detailed Case Facts}

Congress created the food stamp program in 1964 to ``alleviate hunger and malnutrition'' among low-income households and to strengthen the agricultural economy. An eligible ``household'' was originally ``a group of related or non-related individuals'' living as one economic unit. In 1971 a conference committee, without hearings, redefined ``household'' to exclude any group containing a person unrelated to any other member, with narrow exceptions. The scant legislative history showed the amendment aimed at ``hippies'' and ``hippie communes.''

In practice the exclusion fell on destitute households doubling up to survive. Jacinta Moreno, a 56-year-old diabetic on \$75 a month in public assistance, lived with Ermina Sanchez and her three children to share expenses and receive care, leaving her \$10 a month for food; both faced disqualification because they were unrelated. Sheilah Hejny's family faced cutoff for sheltering a troubled 20-year-old woman, and Victoria Keppler lost benefits for sharing an apartment near a school for the deaf so her daughter could attend.

\section{Procedural History}

Moreno and others filed a class action against the Department of Agriculture in the District of Columbia. A three-judge district court held the exclusion unconstitutional under the Fifth Amendment's equal protection component, and the Department appealed directly. The Court heard argument on April 23, 1973, and decided the case on June 25, 1973, the same day as its companion due-process case, \textit{USDA v. Murry}.

\section{Judicial Votes}

Affirmed, \textbf{7--2}.
\begin{itemize}
\item \textbf{Majority:} Justice Brennan, joined by Justices Douglas, Stewart, White, Marshall, Blackmun, and Powell.
\item \textbf{Concurring:} Justice Douglas.
\item \textbf{Dissenting:} Justice Rehnquist, joined by Chief Justice Burger.
\end{itemize}

\section{Holding}

The ``unrelated person'' classification violates the equal protection component of the Fifth Amendment even under rational-basis review. It is irrelevant to the Act's nutritional and agricultural purposes; its actual purpose of excluding ``hippies'' is illegitimate; and the Government's after-the-fact fraud rationale fails because the provision is ``wholly without any rational basis.''

\section{Analysis of Opinions}

\textbf{Majority (Brennan).} Under \textit{Bolling v. Sharpe}, a federal classification must be ``rationally related to a legitimate governmental interest.'' Family relationships have nothing to do with raising nutrition or farm demand, and a purpose to harm hippies is flatly illegitimate. (The Government abandoned its ``morality'' defense.) Departing from the extreme deference of \textit{Dandridge v. Williams}, Brennan also rejected the fraud rationale on two grounds. First, the Act already disqualified able-bodied adults who refuse work and made fraud a felony, casting ``considerable doubt'' on the claim that the amendment targeted the same abuses. Second, because a household is defined by shared cooking and purchasing, any savvy abuser could evade the ban by splitting into two households; in practice it excluded only the most desperately poor, who could not rearrange their lives.

\textbf{Douglas (concurring).} He ``could not say that this `unrelated' person provision has no `rational' relation to control of fraud,'' but the provision burdened the freedom of association---``banding together to combat the common foe of hunger''---triggering strict scrutiny, which it failed.

\textbf{Rehnquist (dissenting).} The majority wrote ``a very persuasive congressional committee report''; under \textit{Dandridge}, Congress could rationally subsidize only related units, since family ties offer rough assurance that a household was not formed to collect benefits. A ``blunt instrument'' is a matter for Congress.

\textbf{Later developments.} \textit{Cleburne} (1985), \textit{Romer v. Evans} (1996), and \textit{United States v. Windsor} (2013) built on \textit{Moreno}'s anti-animus rule, and Justice O'Connor in \textit{Lawrence} described a ``more searching form of rational basis review.'' \textit{Trump v. Hawaii} (2018) retreated, upholding a policy not ``inexplicable by anything but animus.''

\section{Examples: Future Applications}

\begin{enumerate}
\item \textbf{Same-Side (Invalid): Voucher Ban for Protesters.} A city disqualifies households with a protest misdemeanant after its sponsor attacks ``radical campus agitators''; violent felons remain eligible, so the fit is inexplicable except by animus.
\item \textbf{Same-Side (Invalid): Tattoo Artists and Utility Relief.} Excluding households with an unrelated tattoo artist after remarks about ``counterculture freeloaders'' fails, since existing income rules address fraud.
\item \textbf{Opposite-Side (Valid): Neutral Zoning Limit.} A limit on unrelated occupants in single-family zones, as in \textit{Village of Belle Terre v. Boraas}, rests on land-use goals without record animus.
\item \textbf{Opposite-Side (Valid): Work Requirement.} Time-limiting food aid for able-bodied adults without dependents who do not work is tied to the program's purposes and survives deferential review.
\item \textbf{Fence-Sitter: Vehicle-Dwelling Ban.} Sanitation and safety are plausible reasons, but disparaging remarks resemble the anti-hippie history; whether a court sees animus or stray comments is uncertain.
\end{enumerate}

\section{Critique}

\textbf{Doctrinal candor.} Critics, including Rehnquist and Douglas, say \textit{Moreno} claims to apply ordinary rational-basis review while covertly using heightened scrutiny; because the Court named no suspect class, lower courts lack rules for choosing between \textit{Beach Communications} deference and \textit{Moreno}'s ``bite.''

\textbf{Proving animus.} Relying on stray floor remarks raises classic problems of legislative motive: how many legislators must harbor animus, and does an independent rationale such as fraud prevention cleanse a law enacted partly from hostility?

\section{Key Quotations}

\begin{quote}
``For if the constitutional conception of `equal protection of the laws' means anything, it must, at the very least, mean that a bare congressional desire to harm a politically unpopular group cannot constitute a legitimate governmental interest.''
\end{quote}
\noindent --- Justice Brennan, for the Court, 413 U.S. at 534.

\begin{quote}
``Thus, in practical operation, the 1971 amendment excludes from participation in the food stamp program, not those persons who are `likely to abuse the program' but, rather, only those persons who are so desperately in need of aid that they cannot even afford to alter their living arrangements so as to retain their eligibility.''
\end{quote}
\noindent --- Justice Brennan, for the Court, 413 U.S. at 538.

\begin{quote}
``The right of association, the right to invite the stranger into one's home is too basic in our constitutional regime to deal with roughshod.''
\end{quote}
\noindent --- Justice Douglas, concurring, 413 U.S. at 543.

\end{document}
